\begingroup
\chapter*{Évaluation : données et probabilités}\markboth{Évaluation : données et probabilités}{Évaluation : données et probabilités}\addcontentsline{toc}{chapter}{Évaluation : données et probabilités}
\renewcommand{\thebmtexo}{DS4.\arabic{bmtexo}}
\renewcommand{\theHbmtexo}{DS4.\arabic{bmtexo}}
\setcounter{bmtexo}{0}
\begin{devoir}[2 h — 20 points]
Calculatrice autorisée. Justifier les réponses et préciser les domaines utiles. Les questions sont indépendantes sauf indication. 
\end{devoir}
\exo[Évolutions]\label{t11s-ds4-e01}\bareme{4 points}
Un prix vaut $80000$ ariary; il subit une baisse de $15\%$ puis une hausse de $10\%$. Calculer prix final, taux global et taux qui ramènerait le prix final à $80000$ ariary.
\begin{corrige}[exercice~\ref{t11s-ds4-e01}]\label{sol-t11s-ds4-e01}
Prix $80000\times0{,}85\times1{,}1=74800$ (1). Coefficient $0{,}935$, baisse globale $6{,}5\%$ (1). Coefficient retour $80000/74800=200/187$, hausse $13/187\simeq6{,}95\%$ (2).
\end{corrige}
\exo[Statistique]\label{t11s-ds4-e02}\bareme{6 points}
Une série est $4,6,6,8,10,10,12,14$. Calculer moyenne, médiane, quartiles, étendue, variance et écart type. Tracer la boîte à moustaches et interpréter l’écart interquartile.
\begin{corrige}[exercice~\ref{t11s-ds4-e02}]\label{sol-t11s-ds4-e02}
Moyenne $8{,}75$, médiane $9$ (1). Quartiles $6,10$ (1), étendue $10$ (1). Somme des carrés $692$, variance $692/8-(8{,}75)^2=9{,}9375$, écart type environ $3{,}152$ (1). Boîte entre $6,10$, trait $9$, moustaches $4,14$ (1). Écart interquartile $4$ : largeur entre les deux quartiles; avec la convention de rang, au moins la moitié des observations se situe dans $[6;10]$ (1).
\end{corrige}
\exo[Comités et probabilités]\label{t11s-ds4-e03}\bareme{6 points}
Parmi $6$ filles et $4$ garçons, on choisit uniformément un comité de trois élèves.
\begin{enumerate}\item Compter les comités et les comités avec exactement deux filles (2 points).
\item Calculer la probabilité d’au moins une fille (2 points).
\item Choisir un président et un secrétaire distincts parmi les dix élèves donne-t-il le même nombre de résultats que choisir deux élèves sans fonction ? Justifier les deux comptes (2 points).\end{enumerate}
\begin{corrige}[exercice~\ref{t11s-ds4-e03}]\label{sol-t11s-ds4-e03}
1. Total $\binom{10}3=120$; exactement deux filles $\binom62\binom41=60$ (1+1).
2. Sans fille $\binom43=4$, donc $1-4/120=29/30$ (compte 1, complément 1).
3. Fonctions distinctes : $10\times9=90$; groupe de deux : $\binom{10}2=45$, car chaque groupe a deux ordres (1+1).
\end{corrige}
\exo[Réunion d’événements]\label{t11s-ds4-e04}\bareme{4 points}
Pour un dé équilibré, $A$ signifie « résultat pair » et $B$ « résultat supérieur ou égal à $4$ ». Calculer les probabilités de $A$, $B$, $A\cap B$ et $A\cup B$, puis vérifier la formule de réunion.
\begin{corrige}[exercice~\ref{t11s-ds4-e04}]\label{sol-t11s-ds4-e04}
$A=\{2,4,6\}$, $B=\{4,5,6\}$ : $1/2,1/2,1/3,2/3$, puis $1/2+1/2-1/3=2/3$ (1 par étape, formule incluse avec la réunion).
\end{corrige}
\endgroup
